\dresumoheader{mieiro}{Manoel Jorge de Lima}{Atenção e Engajamento em Aulas Remotas: O Desafio da Metodologia Expositiva}{2025}{Trabalho de Conclusão de Curso}{Centro Federal de Educação Tecnológica Celso Suckow da Fonseca}{Rio de Janeiro}{Rio de Janeiro}{2025}

Este estudo visa analisar o grau de dispersão e o baixo engajamento de turmas submetidas a aulas expositivas em ambiente de ensino remoto. Para isso, foi desenvolvida uma extensão para o navegadores \textit{WEB}, chamada ESPEON, que monitora ativamente o comportamento dos alunos durante as aulas, registrando dados de navegação e o uso de periféricos de entrada (câmera e microfone). A pesquisa foi realizada com turmas de cinco alunos, segmentadas por área de conhecimento, expostas a aulas expositivas ministradas por docentes do Centro Federal de Educação Tecnológica Celso Suckow da Fonseca (CEFET/RJ), através da plataforma Microsoft Teams. Ao final de cada aula, a ferramenta gerou relatórios com indicadores com métricas indicativas de grau de atenção e interatividade discente. Os dados quantitativos obtidos foram analisados com estatística descritiva e cruzados com questionários qualitativos sobre a aula, de forma a avaliar o desempenho da metodolgoia expositiva em sala de aula virtual, levando em conta tanto o aspecto objetivo (dados) quanto humano (percepções qualitativas).